\documentclass{article}
\usepackage{graphicx}
\usepackage{amsmath}
\usepackage{amssymb}
\usepackage{amsthm}
\usepackage{hyperref}
\usepackage{listings}
\usepackage{xcolor}
\usepackage{booktabs}
\usepackage{geometry}
\geometry{margin=1in}

% Theorem environments
\newtheorem{definition}{Definition}[section]
\newtheorem{theorem}{Theorem}[section]
\newtheorem{proposition}{Proposition}[section]
\newtheorem{corollary}{Corollary}[section]
\newtheorem{remark}{Remark}[section]

% Code listing style for TAGL
\lstdefinelanguage{TAGL}{
  keywords={},
  sensitive=false,
  comment=[l]{--},
  morecomment=[s]{-*}{*-},
  morestring=[b]",
}
\lstdefinestyle{tagl}{
  language=TAGL,
  basicstyle=\ttfamily\small,
  keywordstyle=\color{blue},
  commentstyle=\color{gray}\itshape,
  stringstyle=\color{teal},
  backgroundcolor=\color{gray!8},
  frame=single,
  framesep=4pt,
  rulecolor=\color{gray!40},
  breaklines=true,
  captionpos=b,
}

% Code listing style for C++
\lstdefinestyle{cpp}{
  language=C++,
  basicstyle=\ttfamily\small,
  keywordstyle=\color{blue},
  commentstyle=\color{gray}\itshape,
  stringstyle=\color{teal},
  backgroundcolor=\color{gray!8},
  frame=single,
  framesep=4pt,
  rulecolor=\color{gray!40},
  breaklines=true,
  captionpos=b,
}

\title{\textbf{The Mathematics of tagd}\\
\large A Formal Description of the Mathematical Structures\\
in the tagd Semantic-Relational System}
\author{Isaac N Colley}
\date{April 2026}

\begin{document}

\maketitle

\begin{abstract}
The tagd system was built from engineering instincts about how knowledge should
be structured. This paper makes the mathematics of those instincts explicit.
Every claim here is grounded in the tagd source code and the TAGL language
specification. The core thesis is simple: \emph{information is structural}. A
tag's meaning is not intrinsic --- it is determined entirely by its position in
the tagspace, by its relations to other tags. This is not a metaphor. It is a
precise mathematical claim, and the sections below give it precise mathematical
form, identifying in the source code the structures of set theory, partial
orders, and Alexandrov topology.
\end{abstract}

\tableofcontents
\bigskip

\section{Primitives}

\subsection{Tags}

A \textbf{tag} is the atomic unit of the tagd system. All tokens in TAGL are
UTF-8 labels \cite{readme}. A tag carries two distinct identities:

\begin{itemize}
  \item \textbf{Nominal identity}: a globally unique UTF-8 label.
  \item \textbf{Structural identity}: its position in the tagspace hierarchy,
        given by its rank, defined by its subordinate relation (aka identity
        relation) \cite{spec}.
\end{itemize}

\begin{definition}[Tag Set]
Let $T$ be the set of all tags in a tagspace. Each tag $t \in T$ is identified
by a string $\mathrm{id}(t) \in \Sigma^*$, where $\Sigma^*$ denotes the set of
all finite UTF-8 strings.
\end{definition}

In TAGL, defining a tag looks like:
\begin{lstlisting}[style=tagl, caption={Defining a tag via a subordinate relation}]
>> dog _is_a mammal;
\end{lstlisting}
Here \texttt{dog} is a tag. Its nominal identity is the string \texttt{"dog"}.
Its structural identity is defined by its subordinate relation to
\texttt{mammal}.

\subsection{The Root and the Axiom}

Every tagspace has a distinguished root tag: \texttt{\_entity}. The README
states:
\begin{quote}
\texttt{\_entity} is axiomatic and self-referencing
(i.e.\ \texttt{\_entity \_sub \_entity}). It is the only tag having this
relation. It is the root of a tagspace.
\end{quote}

\begin{definition}[Axiomatic Root]
The tag \texttt{\_entity} is the \textbf{axiomatic root} of every tagspace. It
is the unique fixed point of the subordinate relation: the one tag whose
super-object is itself.
\end{definition}

In TAGL notation:
\begin{lstlisting}[style=tagl, caption={The axiomatic root}]
_entity _sub _entity
\end{lstlisting}

In set-theoretic terms, \texttt{\_entity} is the \textbf{least element} of the
tagspace under the subordinate order --- the tag that is an ancestor of
everything, including itself.

\section{The Subordinate Relation and the Rooted Tree}

\subsection{Subordinate Relations}

A tag's structural identity is defined by exactly one subordinate (sub)
relation. The TAGL specification states \cite{spec}:
\begin{quote}
A subject MUST have exactly one subordinate relation to exist in the tagspace.
\end{quote}

The sub relators are defined in \texttt{hard-tags.h} \cite{hardtags}:
\begin{lstlisting}[style=cpp, caption={Sub relator definitions in hard-tags.h}]
#define HARD_TAG_SUB     "_sub"     // most abstract sub relator
#define HARD_TAG_IS_A    "_is_a"    // instance subordination
#define HARD_TAG_TYPE_OF "_type_of" // type subordination
\end{lstlisting}

A typical TAGL knowledge hierarchy:
\begin{lstlisting}[style=tagl, caption={Building a subordinate hierarchy in TAGL}]
>> physical_object -^ _entity;
>> living_thing _type_of physical_object;
>> animal _type_of living_thing;
>> mammal _type_of animal;
>> dog _is_a mammal;
\end{lstlisting}

\subsection{The Tagspace as a Rooted Tree}

The subordinate relation organizes all tags into a \textbf{rooted tree}.

\begin{definition}[Tagspace Tree]
A tagspace is a pair $(T, p)$ where:
\begin{itemize}
  \item $T$ is a set of tags containing the distinguished root
        \texttt{\_entity},
  \item $p : T \setminus \{\texttt{\_entity}\} \to T$ is the \emph{parent
        function}, assigning each non-root tag its super-object.
\end{itemize}
There are no cycles (except the axiomatic self-reference of
\texttt{\_entity}).
\end{definition}

The tree structure of the example above:
\begin{verbatim}
_entity
  physical_object
    living_thing
      animal
        mammal
          dog
\end{verbatim}

This is not merely a data structure choice. The tree encodes \emph{taxonomic
knowledge}: to be a \texttt{dog} is to be a \texttt{mammal}, which is to be an
\texttt{animal}, and so on. Structural position \emph{is} meaning.

\section{Rank: Encoding the Tree in Strings}

\subsection{The Rank Function}

Each tag is assigned a \textbf{rank} --- a UTF-8 byte string encoding its
position in the tagspace tree. The \texttt{rank.h} source states:
\begin{quote}
A rank string is the ordinal value of each node in the tagspace tree. A code
point at a particular level is known as the rank value.
\end{quote}

\begin{definition}[Rank Function]
The rank function $r : T \to \Sigma^*$ assigns to each tag a UTF-8 string
such that the string's structure reflects the tag's position in the tree.
\end{definition}

Concrete rank assignments from the bootstrap dump \cite{bootstrap}:

\begin{center}
\begin{tabular}{lll}
\toprule
\textbf{Tag} & \textbf{Rank (dotted)} & \textbf{Position} \\
\midrule
\texttt{\_entity}       & (empty)  & level 0, root \\
\texttt{\_sub}          & 1        & level 1 \\
\texttt{\_is\_a}        & 1.1      & level 2 \\
\texttt{\_type\_of}     & 1.2      & level 2 \\
\texttt{\_rel}          & 2        & level 1 \\
\texttt{\_has}          & 2.1      & level 2 \\
\texttt{\_number}       & 3        & level 1 \\
\texttt{\_integer}      & 3.1      & level 2 \\
\texttt{\_float}        & 3.2      & level 2 \\
\texttt{\_interrogator} & 4        & level 1 \\
\bottomrule
\end{tabular}
\end{center}

\subsection{The Prefix Partial Order}

The key property of ranks is stated explicitly in \texttt{rank.h}:
\begin{quote}
All children of a particular node in the rank hierarchy will be prefixed by the
parent node --- that is, rank substrings define subtrees in the hierarchy.
\end{quote}

\begin{definition}[Containership Relation]
Define the \textbf{containership relation} $\sqsubseteq$ on $T$ by:
\[
  a \sqsubseteq b
  \iff
  r(b) \text{ is a prefix of } r(a).
\]
That is, $a \sqsubseteq b$ means $b$ is an ancestor of $a$ (or $b = a$). Tag
$b$ \emph{contains} tag $a$ in the hierarchy.
\end{definition}

This relation is implemented directly in \texttt{rank.h}:
\begin{lstlisting}[style=cpp, caption={Containership via rank prefix in rank.h}]
// this rank contains the rank being compared,
// in other words, this rank is a prefix to the other
bool contains(const rank&) const;
\end{lstlisting}

And used in \texttt{tagd.cc} for structural operations on tag sets:
\begin{lstlisting}[style=cpp, caption={Containership used in merge\_containing\_tags (tagd.cc)}]
if (b->rank().contains(a->rank())) { // rank can contain itself (b == a)
    merged++;
} else if (a->rank().contains(b->rank())) {
    A.insert(*b);
}
\end{lstlisting}

\subsection{Verifying the Partial Order Axioms}

\begin{theorem}
The containership relation $\sqsubseteq$ is a partial order on $T$.
\end{theorem}

\begin{proof}
We verify the three required properties.

\textbf{Reflexivity.} \texttt{rank.h} notes that \texttt{contains()} holds when
$b = a$, since every string is a prefix of itself. Hence $a \sqsubseteq a$
for all $a \in T$.

\textbf{Antisymmetry.} Suppose $a \sqsubseteq b$ and $b \sqsubseteq a$. Then
$r(b)$ is a prefix of $r(a)$ and $r(a)$ is a prefix of $r(b)$. A string $u$
is a prefix of $v$ and $v$ is a prefix of $u$ if and only if $u = v$. Hence
$r(a) = r(b)$, which implies $a = b$ since ranks uniquely identify tags.

\textbf{Transitivity.} Suppose $a \sqsubseteq b$ and $b \sqsubseteq c$. Then
$r(b)$ is a prefix of $r(a)$ and $r(c)$ is a prefix of $r(b)$. Since the
prefix relation on strings is transitive, $r(c)$ is a prefix of $r(a)$, so
$a \sqsubseteq c$.
\end{proof}

Therefore $(T, \sqsubseteq)$ is a \textbf{poset} (partially ordered set).

\subsection{Lexicographic Total Order}

Ranks also carry a total \textbf{lexicographic order}, implemented in
\texttt{rank.h}:
\begin{lstlisting}[style=cpp, caption={Lexicographic ordering of ranks}]
bool operator<(const rank& rhs) const { return _data < rhs._data; }
\end{lstlisting}

This order --- a deliberate UTF-8 design choice noted in \texttt{rank.h}
(``it collates lexicographically'') --- enables efficient storage and retrieval
of the entire tagspace as a sorted structure. Tags within the same subtree sort
together. The tree \emph{is} the index.

\section{The Tagspace Topology}

\subsection{From Partial Order to Topology}

Every partial order canonically generates a topology. Given the poset
$(T, \sqsubseteq)$, we define:

\begin{definition}[Tagspace Open Sets]
A subset $U \subseteq T$ is \textbf{open} if it is upward-closed under
$\sqsubseteq$:
\[
  U \in \mathcal{T}
  \iff
  \forall\, a \in U,\; \forall\, b \in T :\;
  a \sqsubseteq b \Rightarrow b \in U.
\]
\end{definition}

In the tagspace tree, ``upward-closed'' means: if a tag is in the set, all
its ancestors are too. Each basic open set is the \textbf{subtree rooted at a
tag} together with all tags above it.

This construction is known as the \textbf{Alexandrov topology} --- the topology
generated by a partial order in which arbitrary intersections (not just finite
ones) of open sets remain open.

\subsection{Verifying the Topology Axioms}

\begin{theorem}
Let $M$ be the set of all tags in a tagspace and let $\mathcal{T}$ be the
collection of upward-closed subsets under $\sqsubseteq$. Then
$(M, \mathcal{T})$ is a topological space.
\end{theorem}

\begin{proof}
We verify the three axioms.

\textbf{(i) $\emptyset \in \mathcal{T}$ and $M \in \mathcal{T}$.}
The empty set is vacuously upward-closed. The full set $M$ is upward-closed
because there is no $b \in T$ outside $M$. Both belong to $\mathcal{T}$.

\textbf{(ii) Finite intersections.}
Let $U, V \in \mathcal{T}$. Suppose $a \in U \cap V$ and $a \sqsubseteq b$.
Then $b \in U$ (since $U$ is upward-closed) and $b \in V$ (since $V$ is
upward-closed), so $b \in U \cap V$. Hence $U \cap V \in \mathcal{T}$.

\textbf{(iii) Arbitrary unions.}
Let $\{U_\alpha\}$ be any collection of sets in $\mathcal{T}$. Suppose
$a \in \bigcup_\alpha U_\alpha$ and $a \sqsubseteq b$. Then $a \in U_\alpha$
for some $\alpha$. Since $U_\alpha$ is upward-closed, $b \in U_\alpha
\subseteq \bigcup_\alpha U_\alpha$. Hence $\bigcup_\alpha U_\alpha \in
\mathcal{T}$.
\end{proof}

Therefore $(M, \mathcal{T})$ is a \textbf{topological space}. The
\texttt{merge\_containing\_tags} function in \texttt{tagd.cc} computes
exactly the intersection of two subtrees, exhibiting axiom (ii) in code:
\begin{lstlisting}[style=cpp, caption={Intersection of subtrees in tagd.cc}]
if (b->rank().contains(a->rank())) {
    merged++;
} else if (a->rank().contains(b->rank())) {
    A.insert(*b);
}
\end{lstlisting}

\subsection{Specialization Order}

In the Alexandrov topology, every topological space carries a natural
\textbf{specialization preorder}: $a$ specializes $b$ (written
$a \rightsquigarrow b$) when $a$ is in the topological closure of $\{b\}$.

\begin{proposition}
In the tagspace topology, the specialization order coincides with the
subordinate order: $a \rightsquigarrow b$ if and only if $a$ is a descendant
of $b$ in the tagspace tree.
\end{proposition}

In plain language: \texttt{dog} specializes \texttt{mammal}, which specializes
\texttt{animal}. More specific tags (deeper in the tree) specialize more general
ones. Taxonomic inheritance is not a rule to be enforced --- it is a geometric
consequence of the shape of the space.

\subsection{Continuity as Meaning-Preservation}

\begin{definition}[Tagspace Map]
A map $f : T_1 \to T_2$ between two tagspaces is \textbf{topologically
continuous} if it preserves open sets: $f^{-1}(U) \in \mathcal{T}_1$ for
every $U \in \mathcal{T}_2$.
\end{definition}

\begin{proposition}
A map $f : T_1 \to T_2$ is continuous (with respect to the Alexandrov
topologies) if and only if it is order-preserving: whenever
$a \sqsubseteq b$ in $T_1$, we have $f(a) \sqsubseteq f(b)$ in $T_2$.
\end{proposition}

This gives a precise mathematical definition of a
\textbf{meaning-preserving transformation} between tagspaces. If
\texttt{dog \_is\_a mammal} in $T_1$, a continuous map must send \texttt{dog}
somewhere below the image of \texttt{mammal} in $T_2$.

\section{The Hard Tag Base: A Universal Subspace}

\subsection{Hard Tags}

Every tagspace --- no matter how different its user-defined tags --- contains
exactly the same set of hard tags, defined in \texttt{hard-tags.h}
\cite{hardtags}:

\begin{lstlisting}[style=cpp, caption={Core hard tag definitions in hard-tags.h}]
#define HARD_TAG_ENTITY       "_entity"       // root
#define HARD_TAG_SUB          "_sub"          // subordination
#define HARD_TAG_IS_A         "_is_a"         // instance relation
#define HARD_TAG_TYPE_OF      "_type_of"      // type relation
#define HARD_TAG_RELATOR      "_rel"          // predicate root
#define HARD_TAG_HAS          "_has"          // possession
#define HARD_TAG_CAN          "_can"          // capability
#define HARD_TAG_NUMBER       "_number"       // numeric type root
#define HARD_TAG_INTEGER      "_integer"      // integer type
#define HARD_TAG_FLOAT        "_float"        // float type
#define HARD_TAG_INTERROGATOR "_interrogator" // query root
\end{lstlisting}

\subsection{The Hard Tag Subspace}

\begin{definition}[Hard Tag Subspace]
Let $H \subset T$ denote the set of hard tags in any tagspace $T$. The
induced subspace topology on $H$ is identical in every tagspace. $H$ is
therefore the \textbf{universal invariant subspace} --- a fixed topological
base that every tagspace shares by construction.
\end{definition}

\begin{remark}
In categorical terms, the hard tag subspace is the \textbf{initial object} in
the category of tagspaces: every tagspace is an extension of $H$, and every
valid map between tagspaces must fix $H$.
\end{remark}

\subsection{Interoperability via the Common Base}

Because every tagspace contains $H$ as a common topological foundation, any
two tagspaces can communicate. The hard tags are the shared language --- the
invariant ground that makes the \texttt{httagd} web protocol possible. Two
tagspaces may diverge arbitrarily in their user-defined tags, but they always
agree on the structure of $H$.

\begin{lstlisting}[style=tagl, caption={Hard tags are invariant; user tags vary}]
-- invariant in every tagspace:
_entity _sub _entity

-- user-defined, varies by tagspace:
>> dog _is_a mammal;
\end{lstlisting}

\section{Predicates: Fibered Structure}

\subsection{Relators Are Tags}

Crucially, relators are themselves tags --- subordinate to \texttt{\_rel} in
the tagspace tree. This means predicates are not outside the system; they are
\emph{inside} it. The tagspace is self-describing.

From the bootstrap dump \cite{bootstrap}:
\begin{verbatim}
_rel     rank: 2        (root of all relators)
_has     rank: 2.1
_can     rank: 2.2
_refers  rank: 2.3
\end{verbatim}

The language used to attach meaning to tags is itself expressed in the same
structural language.

\subsection{Predicates as Horizontal Relations}

Predicates are relations attached to a tag that do not alter its position in
the tree. From \texttt{TAGL-spec.md} \cite{spec}:
\begin{quote}
A predicate consists of a relator followed by an object\_list. A relator MUST
be a tag subordinate to \texttt{\_rel}.
\end{quote}

\begin{lstlisting}[style=tagl, caption={Vertical identity and horizontal predicates}]
>> dog _is_a mammal    -- vertical: places dog in the tree
    _has legs = 4, tail
    _can bark;         -- horizontal: predicate data at dog
\end{lstlisting}

\subsection{The Two-Layer Structure}

The tagspace carries two distinct mathematical layers:

\begin{definition}[Two-Layer Tagspace Structure]
\begin{enumerate}
  \item \textbf{Base space}: the topological space $(M, \mathcal{T})$ defined
        by subordinate relations and ranks --- the vertical dimension.
  \item \textbf{Fiber}: at each tag $t \in M$, a set of predicates
        $P(t) \subseteq \mathrm{Rel} \times T \times \mathrm{Mod}$ --- the
        horizontal dimension, where $\mathrm{Rel}$ is the set of relator tags
        and $\mathrm{Mod}$ is the set of modifier values.
\end{enumerate}
\end{definition}

This two-layer structure --- a base space with additional data attached at each
point --- is analogous to a \textbf{fiber bundle} in differential geometry. The
topology gives the skeleton; the predicates give the flesh.

The \texttt{abstract\_tag} class in \texttt{tagd.h} encodes exactly this
decomposition:
\begin{lstlisting}[style=cpp, caption={Two-layer structure in abstract\_tag (tagd.h)}]
class abstract_tag {
    tagd::rank    _rank;       // position in the base space
    predicate_set relations;   // fiber: additional data at this point
};
\end{lstlisting}

\section{Structural Identity}

\subsection{Identity Is Position}

The deepest mathematical claim of the tagd system is that a tag's identity is
entirely structural. The TAGL specification defines:
\begin{quote}
\textbf{Structural identity}: the identity of a tag defined by its subordinate
relation (rank) and predicate\_list.
\end{quote}

A tag \texttt{dog} is not \texttt{dog} because of the string \texttt{"dog"}.
It is \texttt{dog} because of where it sits --- under \texttt{mammal}, under
\texttt{animal}, under \texttt{living\_thing}, under \texttt{\_entity} --- and
because of what predicates are attached to it there. Change the position (the
subordinate relation), and you have a different tag, regardless of the label.

This is \textbf{mathematical structuralism} applied to knowledge
representation: identity is relational, not intrinsic.

\subsection{The Canonical Form}

A tagspace can be fully serialized as TAGL text via \texttt{dump()}. From
the TAGL Thesis document:
\begin{quote}
The canonical form of a tagspace is its dump --- the UTF-8 TAGL code that
represents the tagspace (and can be loaded).
\end{quote}

The dump is not a record of a structure. It \emph{is} the structure, expressed
in TAGL. The tagspace is its own specification.

\section{Summary}

Table~\ref{tab:summary} collects the mathematical structures identified in the
tagd source and maps each to its realization in code and TAGL.

\begin{table}[h]
\centering
\begin{tabular}{lll}
\toprule
\textbf{Mathematical Structure} & \textbf{tagd Realization} & \textbf{Source} \\
\midrule
Set & All tags $T$ in a tagspace & \texttt{tagd.h: tag\_set} \\
Rooted tree & Subordinate relations from \texttt{\_entity} & \texttt{tagd-README.md} \\
Partial order (poset) & Rank prefix containership $\sqsubseteq$ & \texttt{rank.h: contains()} \\
Lexicographic total order & Rank byte string comparison & \texttt{rank.h: operator<} \\
Alexandrov topology & Upward-closed subtrees as open sets & \texttt{rank.h}, \texttt{tagd.cc} \\
Universal base subspace & Hard tag set $H$ & \texttt{hard-tags.h} \\
Fiber / attached structure & Predicate sets $P(t)$ at each tag & \texttt{tagd.h: predicate\_set} \\
Continuous map & Order-preserving tagspace map & \texttt{tagd.cc: merge\_containing\_tags} \\
Canonical serialization & \texttt{dump()} in TAGL canonical form & \texttt{TAGL-spec.md} \\
\bottomrule
\end{tabular}
\caption{Mathematical structures and their tagd source realizations.}
\label{tab:summary}
\end{table}

\section{Conclusion}

The tagd system is not merely a database or a knowledge graph. It is a
\textbf{topological space of structured meaning}. Tags are points. Subordinate
relations give the space its shape. Ranks encode that shape as a partial order.
The Alexandrov topology makes the shape mathematically precise. Predicates
attach additional structure at each point as a fiber over the base topology.
And the hard tag base ensures that every tagspace shares the same foundational
geometry.

TAGL is the language for acting on this space: defining points, querying
regions, transforming structure. Every TAGL statement is a topological
operation.

The vision that \emph{information is structural} --- that meaning comes from
shape, not from labels --- is not a design philosophy. It is a theorem about
how the system works.

\begin{thebibliography}{9}

\bibitem{hardtags}
tagd source code, \texttt{tagd/include/tagd/hard-tags.h}.
Hard tag definitions and axioms.

\bibitem{rankh}
tagd source code, \texttt{tagd/include/tagd/rank.h}.
Rank structure, prefix order, and \texttt{contains()} implementation.

\bibitem{tagdh}
tagd source code, \texttt{tagd/include/tagd.h}.
\texttt{abstract\_tag}, \texttt{predicate\_set}, and tag set operations.

\bibitem{tagdcc}
tagd source code, \texttt{tagd/src/tagd.cc}.
\texttt{merge\_containing\_tags} and merge operations.

\bibitem{spec}
\textit{TAGL Specification}, \texttt{TAGL-spec.md}.
Language specification and semantic model.

\bibitem{readme}
\textit{tagd README}, \texttt{tagd-README.md}.
Working examples and TAGL tutorial.

\bibitem{bootstrap}
\textit{Bootstrap Dump Grid}, \texttt{bootstrap-dump-grid.txt}.
Concrete tagspace rank assignments.

\end{thebibliography}

\end{document}
